\section{行列式的计算规则与推论}
\label{sec:03-Linear-Algebra/04-Determinants/02}

拟合一个高斯过程，每调一次核函数的超参数就要重算对数似然，而似然里坐着一项 \(-\tfrac12\log\det K\)，\(K\) 是 \(2000\times2000\) 的协方差矩阵。若按定义把它展开成排列求和，

\[
\det A=\sum_{\sigma\in S_n}\operatorname{sgn}(\sigma)\prod_{i=1}^n a_{i,\sigma(i)},
\]

就要累加 \(2000!\) 项——这个数大约有五千七百位，把可观测宇宙里的原子逐个数一遍都够不上它的零头。公式精确，也彻底不可用。

出路在上一节的几何里。行列式是体积因子，而消元的每一步都对应一个看得懂的几何动作，因此只要一边消元一边把每步的影响记下来，行列式会自己走到终点。

\subsection{消元把行列式一路记下来}

取

\[
A=\begin{pmatrix}1&2&0\\2&5&1\\1&0&3\end{pmatrix},
\]

依次做 \(R_2\leftarrow R_2-2R_1\)、\(R_3\leftarrow R_3-R_1\)、\(R_3\leftarrow R_3+2R_2\)。三步都是行倍加，也就是纯剪切，行列式一动不动。终点是上三角矩阵

\[
U=\begin{pmatrix}1&2&0\\0&1&1\\0&0&5\end{pmatrix},
\qquad
\det A=\det U=1\cdot1\cdot5=5.
\]

记账规则很短：过程中每交换一次行，末尾补一个负号；每把某行整体乘过 \(c\)，末尾除以一个 \(c\)。代价是 \(O(n^3)\) 次运算，和解一个线性方程组同一个量级，比 \(n!\) 好得不像话。

如果消元结果已经以 \(PA=LU\) 的形式存下来了，那连账都不用重记。\(L\) 是单位下三角，对角元全为 \(1\)，行列式为 \(1\)；置换矩阵 \(P\) 的行列式是 \(\pm1\)，由交换次数的奇偶决定。于是 \(\det A=\pm\prod_i u_{ii}\)。换句话说，只要做过一次分解，行列式就是免费的副产品。

三角矩阵的行列式等于对角元之积，上一节从``剪切不改变体积''读出来过；从排列公式也能看出来，因为非零项要求每个 \(\sigma(i)\ge i\)，唯一满足的排列是恒等排列。两条路给出同一个结论。

\subsection{乘法定理是一句几何话}

\begin{center}
\begin{tikzpicture}[x=1cm,y=1cm,font=\small,>=stealth]
  \begin{scope}
    \draw[black!55,fill=blue!10] (0,0) -- (1,0) -- (1,1) -- (0,1) -- cycle;
    \node[font=\footnotesize] at (0.5,-0.6) {面积 \(1\)};
  \end{scope}
  \draw[->,black!70] (1.5,0.8) -- (2.6,0.8) node[midway,above,font=\footnotesize] {\(B\)};
  \begin{scope}[shift={(3.1,0)}]
    \draw[black!25,dashed] (0,0) -- (1,0) -- (1,1) -- (0,1) -- cycle;
    \draw[black!55,fill=blue!14] (0,0) -- (1.2,0.3) -- (1.6,1.3) -- (0.4,1.0) -- cycle;
    \node[font=\footnotesize] at (0.8,-0.6) {面积 \(\abs{\det B}\)};
  \end{scope}
  \draw[->,black!70] (5.3,0.8) -- (6.4,0.8) node[midway,above,font=\footnotesize] {\(A\)};
  \begin{scope}[shift={(6.9,0)}]
    \draw[black!25,dashed] (0,0) -- (1.2,0.3) -- (1.6,1.3) -- (0.4,1.0) -- cycle;
    \draw[black!55,fill=blue!20] (0,0) -- (1.38,0.96) -- (2.38,2.36) -- (1.0,1.4) -- cycle;
    \node[font=\footnotesize] at (1.2,-0.6) {面积 \(\abs{\det A}\abs{\det B}\)};
  \end{scope}
\end{tikzpicture}

\emph{\(B\) 先把单位正方形的面积乘上 \(|\det B|\)；\(A\) 接手时面对的已经是一个歪斜的平行四边形，但上一节那句``所有区域共享同一个因子''保证它照样只乘 \(|\det A|\)。两步下来，复合变换 \(AB\) 的因子就是两个因子的乘积。}
\end{center}

于是

\[
\det(AB)=\det A\,\det B.
\]

定向的账也一样：翻两次面等于没翻，符号相乘。这个论证的关键一步是``对任意区域用同一个因子''，若只知道 \(A\) 对单位立方体的作用，就接不上 \(B\) 输出的那个歪斜形状。

严格的代数证明走另一条路：先对三类初等矩阵逐一验证乘法性质，再把可逆矩阵写成初等矩阵之积；\(A\) 或 \(B\) 奇异时两边都是零，等式仍然成立。条件对账：这条等式只要求 \(A,B\) 是同阶方阵，可逆与否无所谓。但它对加法完全失效——取 \(A=B=I_2\)，左边 \(\det(2I_2)=4\)，右边 \(\det A+\det B=2\)。行列式对单独某一列线性，对整个矩阵并不线性，这两件事不要混。

乘法定理一旦到手，几条常用推论就是把它用一遍的事。对 \(AA^{-1}=I\) 取行列式，立刻得到 \(\det(A^{-1})=1/\det A\)：把体积撑大三倍的变换，撤销时必然把体积压回三分之一。连用两次则有 \(\det(S^{-1}AS)=\det A\)，换基不改变行列式——坐标变换先把东西拉过去，末了又原样拉回来，中间那个体积因子是变换自身的属性，与用哪组基描述它无关。这也预告了下一单元的一个等式：行列式等于全部特征值之积，因为特征值同样不随换基而变。

还有一条 \(\det(A^{\trans})=\det A\)，它不经过乘法定理。在排列公式里，\(A^{\trans}\) 中由 \(\sigma\) 贡献的项与 \(A\) 中由 \(\sigma^{-1}\) 贡献的项逐个相同，而 \(\operatorname{sgn}(\sigma^{-1})=\operatorname{sgn}(\sigma)\)，两个和式重排后完全一致。这条等式补上了上一节欠下的一笔账：列操作的规则与行操作相同，因为列操作就是转置之后的行操作。

\subsection{余子式展开：知道它是什么就够了}

固定第 \(i\) 行，用多线性把它拆成 \(n\) 个只在一个位置非零的行。对只含 \(a_{ij}\) 的那一项，用交错性把第 \(j\) 列与第 \(i\) 行逐次挪到左上角，共 \(i+j-2\) 次相邻交换，符号 \((-1)^{i+j}\) 由此而来；余下的体积由划掉第 \(i\) 行第 \(j\) 列所得的子矩阵 \(M_{ij}\) 承担。于是

\[
\det A=\sum_{j=1}^n a_{ij}\,(-1)^{i+j}\det M_{ij}.
\]

这就是 Laplace 展开，沿任意一行或任意一列都成立，手算时挑零最多的那一行最省事。

作为算法它没有前途。展开一次要算 \(n\) 个 \(n-1\) 阶行列式，递归下去总项数恰好是 \(n!\)，与排列公式一样糟。它的价值在别处：四阶以内手算、带符号参数的小矩阵、以及推出伴随矩阵的恒等式 \(A\operatorname{adj}(A)=(\det A)I\)，进而在 \(\det A\neq0\) 时给出

\[
A^{-1}=\frac{1}{\det A}\operatorname{adj}(A).
\]

这个封闭式在证明里好用，比如它显示逆矩阵的每个元素都是原矩阵元素的有理函数，因而对参数光滑依赖。但没有任何数值软件按它求逆。

Cramer 法则的处境相同。可逆时 \(Ax=b\) 的第 \(j\) 个分量等于

\[
x_j=\frac{\det A_j(b)}{\det A},
\]

其中 \(A_j(b)\) 是把 \(A\) 的第 \(j\) 列换成 \(b\) 所得的矩阵。推导只用一次多线性：把 \(b=\sum_k x_ka_k\) 代进去展开，\(k\neq j\) 的项都因出现两个相同的列而消失，只剩 \(x_j\det A\)。公式很好地说明了 \(\det A=0\) 时唯一解表达式为什么失效。可它要算 \(n+1\) 个行列式，成本远高于一次分解加前代回代，实际求解从不这么做，理由见 \hyperref[sec:08-Numerical-Analysis/07-Numerical-Methods-for-ML/03]{``需要的是解，而不是逆矩阵''}。

\subsection{体积因子在密度变换里现身}

线性替换 \(x=Au\) 把体积元乘上 \(|\det A|\)，写作 \(dx=|\det A|\,du\)。非线性映射 \(x=g(u)\) 在每一点附近被 Jacobian 矩阵 \(J_g(u)\) 线性近似，局部体积因子就是 \(|\det J_g(u)|\)——只是这个因子不再是常数，它随位置而变。多重积分的换元公式因此带上这一项，完整讨论见 \hyperref[sec:02-Calculus/06-Multivariable-Calculus/07]{``多重积分与变量替换''}。

\begin{center}
\begin{tikzpicture}[x=1cm,y=1cm,font=\small,>=stealth]
  \begin{scope}
    \draw[black!35,step=0.4] (0,0) grid (1.6,1.6);
    \draw[black!60,fill=blue!16] (0.8,0.8) rectangle (1.2,1.2);
    \node[font=\footnotesize] at (0.8,-0.5) {\(u\) 平面};
  \end{scope}
  \draw[->,black!70] (2.2,0.8) -- (3.3,0.8) node[midway,above,font=\footnotesize] {\(g\)};
  \begin{scope}[shift={(3.9,0)}]
    \foreach \c in {0,0.4,0.8,1.2,1.6}{
      \draw[black!35] plot[domain=0:1.6,samples=25] ({\c+0.25*\x*\x},{\x+0.2*\c*\c});
      \draw[black!35] plot[domain=0:1.6,samples=25] ({\x+0.25*\c*\c},{\c+0.2*\x*\x});
    }
    \draw[black!60,fill=blue!16] (0.96,0.928) -- (1.36,1.088) -- (1.56,1.488) -- (1.16,1.328) -- cycle;
    \node[font=\footnotesize] at (1.2,-0.5) {\(x=g(u)\) 平面};
  \end{scope}
\end{tikzpicture}

\emph{均匀的方格网被非线性映射拧成曲线网。每个小方格近似变成一个平行四边形，面积乘以 \(|\det J_g(u)|\)；网格越密，近似越准。这里的因子随位置递减：高亮的那一格面积只剩原来的 \(0.8\) 倍，而原点附近几乎没有变化。}
\end{center}

概率密度用的是这条规则的倒数形式。设 \(u\) 的密度为 \(p_u\)，\(g\) 可逆，则

\[
p_x(x)=p_u\bigl(g^{-1}(x)\bigr)\,\bigl|\det J_g\bigl(g^{-1}(x)\bigr)\bigr|^{-1}.
\]

概率总量守恒，空间被撑开的地方密度只能摊薄，被压紧的地方密度只能升高。若某处 \(\det J_g=0\)，映射在那里把一个方向压平，密度会发散——这也是可逆性在换元公式里不是可有可无的原因。

生成模型中的 normalizing flow 把这条式子当成架构约束来用。模型是一串可逆变换的复合，训练目标是数据的对数似然，里面含有 \(\sum_k\log\abs{\det J_k}\)；乘法定理保证复合映射的 Jacobian 行列式等于各层之积，取对数就变成求和。麻烦在于一般的 \(n\times n\) 行列式要 \(O(n^3)\)，每个样本每一层都算一次根本负担不起。于是这类模型的设计几乎都绕着同一件事转：让每层的 Jacobian 是三角的或对角的，行列式退化为对角元连乘，代价降到 \(O(n)\)。耦合层、自回归流、以及把权重矩阵参数化成 \(PLU\) 形式的可逆 \(1\times1\) 卷积，动机都在这里。

取对数还有一个纯数值的理由。高维矩阵的行列式动辄溢出或下溢：\(1000\) 个 \(0.1\) 相乘就是 \(10^{-1000}\)，双精度浮点根本存不下，算出来只会得到零。所以多元正态的对数密度里写的是 \(\log\det\Sigma\)，而实现时由 Cholesky 分解 \(\Sigma=LL^{\trans}\) 给出

\[
\log\det\Sigma=2\sum_{i=1}^n\log \ell_{ii},
\]

全程在对数尺度上做加法，不会溢出，也顺带保证了 \(\Sigma\) 正定（见 \hyperref[sec:03-Linear-Algebra/06-Symmetric-and-Positive-Definite-Matrices/03]{``Cholesky 分解：把正能量写成平方和''}）。\(\log\det\Sigma\) 本身也可以直接读：它衡量协方差椭球有多大，也就是这个分布把不确定性摊在了多大一块空间上。实验设计里的 D 最优准则要最大化 Fisher 信息矩阵的行列式，动机正是把参数的置信椭球压到体积最小。

行列式给出的始终是一个总量：所有方向的缩放乘在一起。判断可逆、校正密度、度量椭球大小，它都够用；可它把方向信息抹平了，哪个方向被拉长、哪个方向快要塌掉，一个标量装不下。下一单元 \hyperref[ch:021]{``特征值、对角化与动力系统''}把这个乘积拆开：特征值给出各个不变方向上的缩放倍数，它们的积正是行列式，它们的和是迹。
